\documentclass[main.tex]{subfiles}
\begin{document}
\section{Software}
\begin{myTheorem}{Software}{thm:software}
  Die Programme, Verfahren, zugehörige
Dokumentation und Daten, die mit dem Betrieb
eines Rechnersystems zu tun haben. (IEEE 610.12)
\end{myTheorem}
\subsection*{Typen von Software}
\begin{itemize}
    \item \textbf{Individualsoftware}: 
    \begin{itemize}
        \item Software mit spezifischer Funktionalität, die für Auftraggeber*innen entwickelt wird
    \end{itemize}
    \item \textbf{Softwareprodukte}: 
    \begin{itemize}
        \item Softwareprodukte werden für einen Markt mit einem bestimmten Bedarf entwickelt
        \item Normalerweise keine direkten Auftraggeber*innen
    \end{itemize}
    \item \textbf{Softwareplattform}: 
    \begin{itemize}
        \item Ein Softwareprodukt, dass den Aufbau neuer Anwendungen in derselben Plattform erlaubt
    \end{itemize}
    \item \textbf{Software realisiert Systeme}:
    \begin{itemize}
        \item Mengen von Programmen genügen nicht; wir brauchen Systeme
        \item Zusammenspiel vieler verteilter Bausteine aus Hardware und Software notwendig
        \item Durch die (notwendige) Systemintegration wird aus vielen einfachen kleinen Teilproblemen ein komplexes und schwieriges Gesamtproblem!
    \end{itemize}
\end{itemize}
\subsection*{Eigenschaften von Software}
\begin{itemize}
    \item \textbf{immateriell}: Software ist abstrakt und mit unseren Sinnen nicht direkt wahrnehmbar!
    \begin{itemize}
        \item Fehler sind schwieriger erkennbar
        \item Entwicklungsstand sowie Qualität schwer zu beurteilen
    \end{itemize}
    \item \textbf{unstet}
    \begin{itemize}
        \item Kleinste Veränderungen der Software führen ggf. zu massiven Änderungen im Systemverhalten
        \item Zusammenhang zwischen Eingabe und Ausgabe nicht durch eine kontinuierliche Funktion beschreibbar
        \item Valider Funktionsnachweis schwierig, insbesondere im Nachhinein, ggf. schwieriger als in anderen Domänen
        \item Nicht-triviale Softwaresysteme sind intrinsisch komplex
    \end{itemize}
    \item \textbf{benötigt Qualität}
    \begin{itemize}
        \item \underline{Zuverlässigkeit}: Software darf im Fall des Versagens keine physischen oder ökonomischen Schäden verursachen
        \item \underline{Benutzbarkeit}: Software muss sich nach den Bedürfnissen der Benutzer richten, die Benutzerschnittstelle muss ergonomisch und selbsterklärend sein und die Dokumentation muss zielgerecht zur Verfügung stehen
        \item \underline{Wartbarkeit}: Software muss an neue Anforderungen anpassbar sein und sollte möglichst plattformunabhängig sein
        \item \underline{Effizienz} Software muss ökonomischen Gebrauch von Ressourcen des unterliegenden Systems machen
        \item \underline{nicht-funktionale Qualitätseigenschaften} \ref{sec:vl3}
    \end{itemize}
    \item \textbf{verändert Realität}
    \begin{itemize}
        \item \underline{Automatisierung oder Unterstützung} von menschlicher Arbeit, technischen oder sozialen Abläufen und in beliebigen Anwendungsbereichen
        \item in ständiger \underline{Wechselwirkung} mit Arbeitsprozessen, Produktionsprozessen und Menschen 
        \item hat zwei hauptsächliche \textbf{Herausforderungen}: Das Problem im Kontext des jeweiligen Anwendungsbereichs verstehen und die Problemlösung auf Software abbilden
    \end{itemize}
\end{itemize}
\subsection*{Der Software Lebenszyklus}
\begin{itemize}
    \item \textbf{Analyse und Spezifikation:} Was soll das System machen?
    \item \textbf{Entwicklung:} Erstellung des Systems, Integration und Test.
    \item \textbf{Abnahme:} Übereinstimmung mit den Kundenanforderungen.
    \item \textbf{Betrieb:} Installation beim Kunden, Wartung.
    \item \textbf{Evolution:} Weiterentwicklung gemäß neuer Kundenanforderungen oder Ausmusterung.
\end{itemize}

\section{Software Engineering}
\begin{myTheorem}{Software Engineering}{thm:se}
Software Engineering ist die Anwendung eines
\textbf{systematischen}, \textbf{disziplinierten} und \textbf{quantifizierbaren}
Ansatzes auf die \textbf{Entwicklung}, den \textbf{Betrieb} und die \textbf{Wartung}
von Software, das heißt, die Anwendung der Prinzipien des
Ingenieurwesens auf Software. (IEEE 610.12)
\end{myTheorem}

\begin{itemize}
    \item \textbf{Systematisch}: Klar abgegrenzte Einzelschritte und Ein- und Ausgaben sind definiert
    \item \textbf{Diszipliniert}: Vorgehen folgt organisierten Prinzipien, feste Rollendefinitionen, sorgfältige Dokumentation
    \item \textbf{Quantifizierbar}: Schätzung und Messung von relevanten Größen und vorhersagbar mit begründbaren und akzeptablen Grenzen
\end{itemize}

\subsection*{Typen von Software Engineering} 
\begin{figure}
    \centering
    \begin{subfigure}
    \includegraphics{}
    \caption{Caption}
    \label{fig:my_label}
    \end{subfigure}
    \begin{subfigure}
    \includegraphics{}
    \caption{Caption}
    \label{fig:my_label}
    \end{subfigure}
    \begin{subfigure}
    \includegraphics{}
    \caption{Caption}
    \label{fig:my_label}
    \end{subfigure}
    \caption{Caption}
    \label{fig:my_label}
\end{figure}
\begin{table}[H]
    \centering
   \begin{tabular}{p{1.8cm} p{3.8cm} p{3.8cm} p{3.8cm} } 
    \textbf{Merkmal} & \textbf{Projektbasiertes Software Engineering}   & \textbf{Produktorientiertes Software Engineering} & \textbf{Produktmanagement in Softwareprojekten}  \\  \hline
    \scriptsize{Ausgangspunkt} & Kundenanforderungen & Geschäftsidee oder Marktlücke & Interne strategische Planung \\  
    \scriptsize{Zielsetzung} & Individuelle Software, die spezifische Anforderungen erfüllt & Entwicklung eines marktfähigen Softwareprodukts & Steuerung der Entwicklung, Vermarktung und Wartung eines Produkts  \\  
    \scriptsize{Kundenfokus} & Externer Kunde mit spezifischen Anforderungen & Kein externer Kunde; Fokus auf Marktanforderungen & Schnittstelle zwischen Kunden, Unternehmen und Entwicklungsteam \\  
    \scriptsize{Flexibilität} & Anforderungen können sich während der Entwicklung ändern & Time-to-Market hat Priorität, daher weniger flexibel & Anpassung an Kundenbedürfnisse und Marktanforderungen \\  
    \scriptsize{Lebensdauer der Software} & Lange Lebensdauer (10 Jahre oder mehr) mit langfristiger Unterstützung & Abhängig von Marktakzeptanz und Produktstrategie & Überwachung des gesamten Produktlebenszyklus \\ 
    \scriptsize{Entscheidungs- findung} &Kunde gibt Anforderungen vor & Entscheidungen durch internes Produktmanagement & Produktmanager*innen tragen Gesamtverantwortung \\  
    \scriptsize{Hauptfokus} & Kundenorientierte Lösung &  Marktgerechtes Produkt& Langfristige Produktstrategie und -planung \\  
\end{tabular}
    \caption{Typen von Software Engineering}
    \label{tab:setypen}
\end{table}

\subsection*{Softwareentwicklungsprozesse}
\begin{myTheorem}{Softwareprozess}{thm:softwareprozess}
Satz von Aktivitäten dessen Ziel die Entwicklung oder
Weiterentwicklung von Software ist.
\end{myTheorem}
\begin{myTheorem}{Software development process}{thm:softwareprocess}
The process by which user needs are translated
into a software product. The process involves translating user needs into software
requirements, transforming the software requirements into design, implementing
the design in code, testing the code, and sometimes, installing and checking out the
software for operational use. (IEEE 610.12)
\end{myTheorem}
\subsubsection*{Naives Vorgehen}
\begin{table}[H]
    \centering
    \begin{tabular}{p{2cm} p{13.5cm}}
           \scriptsize{Name} & \textbf{Code-and-Fix}  \\ \\
       \scriptsize{Vorgehen}  & Vorläufiges Programm erstellen; Anforderung, Entwurf, Testen, Wartung überlegen; Programm entsprechend verbessern \\ \\
       \scriptsize{Eigenschaften} & \textbf{keine}: Spezifikation, Entwurf, nachvollziehbare Systematik \\  \\
       \scriptsize{Vorteile} & schnell, einfach\\ \\
       \scriptsize{Nachteile} & nicht planbar, Teamarbeit unmöglich, Programme schlecht strukturiert, keine bestimmten Anforderungen, extrem hoher Korrekturaufwand, schlechte Dokumentation - also \textbf{schlechte Qualität} und \textbf{teuer}
    \end{tabular}
\end{table}
\subsubsection*{Generelle Schritte}
\begin{enumerate}
    \item Anwendung
    \begin{itemize}
        \item Nutzung
        \item Sonderfälle
        \item Performance
        \item Technologische oder regulatorische Vorgaben
    \end{itemize}
    \item Struktur
    \begin{itemize}
        \item Benutzeroberfläche
        \item Architektur
        \item Klassen
        \item Algorithmik
        \item Programmiersprache
    \end{itemize}
    \item Programmierung
    \begin{itemize}
        \item Modularität
        \item Nutze Code-Standards und Muster (Erklärbarkeit)
        \item Standardbibliotheken 
        \item Fehlerbehandlungen
    \end{itemize}
    \item Testen 
    \begin{itemize}
        \item Korrekte Funktionalität
        \item Robustheit und Fehlervermeidung
        \item Performance-Aspekte
        \item Erweiterbarkeit
    \end{itemize}
    \item Aufwandsschätzung
    \begin{itemize}
        \item Aufwände, benötigte Ressourcen und Dauer
        \item eine Implementierungszeit und Aspekte wie Pausen, Tests, Debugging und reale Faktoren wie Verzögerungen
    \end{itemize}
\end{enumerate}
\subsubsection*{Beispiel}
\begin{table}[H]
    \centering
    \begin{tabular}{p{2cm} p{13.5cm}}
        \scriptsize{Name} & Wasserfallmodell  \\ \\
        \scriptsize{Typ} & dokumentengetriebenes Modell \\ \\
       \scriptsize{Vorgehen}  & \\ \\
       \scriptsize{Eigenschaften} &  \\  \\
       \scriptsize{Vorteile} &  \\ \\
       \scriptsize{Nachteile} & 
    \end{tabular}
\end{table}
\subsubsection*{Vorgehens- vs. Prozessmodelle}
\begin{table}[H]
    \centering
    \begin{tabular}{p{2cm} p{6.75cm} p{6.75cm}}
        \scriptsize{Name} & \textbf{Vorgehensmodell} & \textbf{Prozessmodell} \\ \\
        \scriptsize{Umfang} & Auszuführende Entwicklungsschritte und deren Reihenfolge  & Verfahren zur Anforderungserhebung und Prüfung von Zwischenergebnissen,  Projektphasen, Meilensteine und Prüfkriterien, Notationen und Terminologie, Werkzeuge\\ \\
        \scriptsize{Fokus} & Praktische Umsetzung der Entwicklung & Umfassende Sicht auf den gesamten Softwareentwicklungsprozess \\ \\
        \scriptsize{Flexibilität} & Oft weniger formal und flexibler in der Anwendung & Strukturierter und formeller Ansatz \\ \\
    \end{tabular}
\end{table}
%\section{Motivation}
%\section{Konzept}
\end{document}